\section{Reviewer B Comments}

{\it
``First, I would like to apologize to the authors for the time it took me to complete this review.
This manuscript investigates a moderately promising approach toward improving an important
aspect of reactor-monitoring safeguards using neutrinos. In particular, the potential use of inverse
beta decay (IBD) directionality may have broader implications for other areas, including
supernova neutrino detection, as briefly discussed by the authors. This is especially relevant
given the limitations of the other existing dominant technologies in reconstructing the
directionality of the neutrino component of the supernova spectrum. The authors provide a clear
and appropriate contextualization of the work, explicitly state the objectives, and identify the
main challenges currently facing the field. I therefore recommend the manuscript for publication
in Physical Review A.
I have only a few minor comments, listed below.

\begin{enumerate}
    \item {\bf Abstract:} \\It would be appropriate to explicitly state that the results presented in this work are obtained for lithium-6–doped detectors. In addition, the authors should clarify that they anticipate the methodology may be generalizable to other IBD detector technologies. However, I am skeptical regarding the degree to which such generalization is valid across alternative detector implementations.
    \item  {\bf Page 2:} \\A more explicit discussion of the potential impact of the limited spatial distribution of lithium-6 doping in SoLid and CHANDLER detectors on the presented results would improve the clarity of the manuscript.
    \item  {\bf Conclusion:} \\It would be beneficial to include a more detailed discussion of the authors’ intended next steps and possible future research directions arising from this work. In particular, while this paper represents a meaningful first step, it would be helpful to outline future directions aimed at improving the mathematical robustness of the method.
\end{enumerate}
''
}

\section{Response}

Dear Reviewer B:

Thank you for your comments and we appreciate your feedback. Rest assured that we understand that the end of semester can be a busy time for reviewers, so we appreciate you getting back to us. We also appreciate the recommendation for publication in PR Applied, and will incorporate your suggestions in our revised manuscript. Please find our responses to your comments in the order presented:

\begin{itemize}
    \item Abstract \\ The abstract has been changed to read: ``We treat data from Lithium-6-doped detector segments in the form of a matrix.'' -- specifying that the simulated detector data are from a lithium-6-doped detector. Our intent to generalize the method to other detector implementations is addressed in the conclusion.
    \item Page 2 \\ The following statement has been added explaining that the limited spatial distribution of the neutron detection layers in SOLID and CHANDLER was not studied during the development of this algorithm: \\
    \begin{quote}
         This confinement of the ${}^6$Li dopant to the surfaces of the segments in these detectors makes the neutron detection layer much thinner. This thin neutron detection layer contradicts the implicit assumption of isotropy in the binning scheme introduced in Figs. \ref{fig_segmentation_scaling} and \ref{fig:phi_def}. Modifying the binning and the directionality algorithm to address these anisotropies is not within the scope of this work.
    \end{quote}
    \item Conclusion \\ The planned studies to apply the method to other detector technologies and search for ways to improve mathematical robustness has been more explicitly stated with the following final sentence: \\
    \begin{quote}
        Planned future studies will investigate incorporating various statistical methods in directional analysis, applying this method to detectors of different geometry such as monolithic and separated, as well as with different detection media such as water Cherenkov, and applying the method to three dimensionally segmented detectors. 
    \end{quote}
\end{itemize}

Sincerely,

Brian Crow
Department of Physics and Astronomy, University of Hawai'i at M\={a}noa

on behalf of the authors